%% Determines the type of document (including standard settings for layout).
\documentclass[letterpaper]{article}



%% Package to control the size of a page: the standard margins used by LaTeX are
%% very wide!
\usepackage[margin=1in]{geometry}


%% Packages add functionality to the document. The AMS packages are standard
%% packages to support various mathematical notations. AMS stands for American
%% Mathematical Society, the organization that maintains these packages.
\usepackage{amsmath,amsthm,amssymb}

%% Theorems-like environments using functionality provided by amsthm.
\theoremstyle{plain}
\newtheorem{theorem}{Theorem}[section]
    %% [section] at the end specifies that theorems should be numbered
    %% per-section: Section x starts with theorem-like x.1 and so on, ...
\newtheorem{proposition}[theorem]{Proposition}
    %% [theorem] in the middle specifies: use the same counter as the theorem
    %% environment: here we number all theorem-like environments consecutively.
\newtheorem{corollary}[theorem]{Corollary}
\newtheorem{lemma}[theorem]{Lemma}
\theoremstyle{definition}
\newtheorem{definition}[theorem]{Definition}
\theoremstyle{remark}
\newtheorem{example}[theorem]{Example}
\newtheorem{remark}[theorem]{Remark}


%% Support for nicely formatted tables.
\usepackage{booktabs}


%% Support for colors & colors in tables.
\usepackage[table]{xcolor}

% Seven colors safe for use color blindness.
% Colors taken from doi:10.1038/nmeth.1618.
\definecolor{cbOrange}{RGB}{230,159,0}
\definecolor{cbSkyBlue}{RGB}{86,180,233}
\definecolor{cbBluishGreen}{RGB}{0,158,115}
\definecolor{cbYellow}{RGB}{240,228,66}
\definecolor{cbBlue}{RGB}{0,114,178}
\definecolor{cbVermillion}{RGB}{213,94,0}
\definecolor{cbReddischPurple}{RGB}{204,121,167}


%% Notation used in this document.
\newcommand{\n}{\mathbf{n}} %% Num. Replicas.
\newcommand{\f}{\mathbf{f}} %% Num. Faulty Replicas.

%% Misc. Math notation.
\newcommand{\BigO}{\mathcal{O}}
\newcommand{\abs}[1]{\lvert #1 \rvert}
\newcommand{\AName}[1]{\textsc{#1}}
\newcommand{\Var}[1]{\texttt{#1}}


%% Algorithms.
\usepackage{algorithm}
\usepackage[noend]{algorithmic}
\newcommand{\GETS}{:=}


%% Formatting SI-units.
\usepackage{siunitx}
\sisetup{per-mode=symbol}


%% TikZ: for creating figures.
\usepackage{tikz}

%% Configuration for figures: Nicer arrows.
\usetikzlibrary{arrows.meta}
\tikzset{>=Stealth}


%% pgfplots: drawing plots using TikZ.
\usepackage{pgfplots}
%% Configuration for plots: Use color-blind friendly colors.
\pgfplotscreateplotcyclelist{cbSafeList}{
    very thick,solid,cbOrange,every mark/.append style={solid},mark=*\\
    very thick,solid,cbSkyBlue,every mark/.append style={solid},mark=*\\
    very thick,solid,cbBluishGreen,every mark/.append style={solid},mark=*\\
    very thick,solid,cbYellow,every mark/.append style={solid},mark=*\\
    very thick,solid,cbBlue,every mark/.append style={solid},mark=*\\
    very thick,solid,cbVermillion,every mark/.append style={solid},mark=*\\
    very thick,solid,cbReddischPurple,every mark/.append style={solid},mark=*\\
    very thick,solid,black,every mark/.append style={solid},mark=*\\
}
\pgfplotsset{
    legend style={font=\small},
    compat=1.16,
    width=260pt,
    height=140pt,
    legend cell align=left,
    xlabel near ticks,
    ylabel near ticks,
    every axis/.append style={
        cycle list name=cbSafeList,
        ymin=0,
        enlargelimits=0.05,
        mark size=1pt,
        ylabel style={align=center},
        xlabel style={align=center},
        title style={align=center}
    }
}

%% PgfplotsTable: loading data files to use with pgfplots.
\usepackage{pgfplotstable}


%% Support for hyperlinks and urls. The setting ``colorlinks'' sets how links
%% are shown in the document (with a color, without underline). We put hyperref
%% last---it has a tendency to break other packages when loaded before them.
\usepackage[colorlinks]{hyperref}
\usepackage{graphicx}
\usepackage{pgffor}
\usepackage{caption}
\usepackage{tabularx}
\usepackage{enumitem}
\usepackage{fancyhdr}

\begin{document}

%% First the meta-data of the doucment: title,

\begin{titlepage}
    \centering
    \vspace*{3cm}
    
    {\Huge COMPSCI 2DB3 Assignment 6}\\[2cm]
    
    {\large Luca Mawyin}\\[0.5cm]
    
    {\large Dr. Jelle Hellings}\\[0.5cm]

    {\large \today}\\[2cm]

    {\Large McMaster University}

    \vfill
    \begin{flushleft}

        {\large \textbf{Student Number: }400531739}\\[0.5cm]
        {\large \textbf{MacID: }mawyinl}
    \end{flushleft}

\end{titlepage}
\renewcommand{\thesection}{Part \arabic{section}}
\renewcommand{\thesubsection}{Problem \arabic{subsection}}
\renewcommand{\thesubsubsection}{P\arabic{subsection}.\arabic{subsubsection}}

\section{}
\subsection{}
\subsubsection{}

\begin{itemize}
    \item Every \textit{rid} has a (\textit{title, total-duration, min-price}). \\[0.5em]
    \textbf{Explanation:} (\textit{title, total-duration, min-price}) are descriptive of the recipe as a whole, and therefore would remain constant for all entries with the same \textit{rid}.
    \item Every recipe has at least one step \\[0.5em]
    \textbf{Explanation:} The description states that a recipe has one or more steps, and therefore cannot exist without at least one step.
    \item Every (\textit{rid, step-order}) determines a unique (\textit{description, duration})\\[0.5em]
    \textbf{Explanation:} Every recipe has a unique sequence of steps, and each step has a description and duration. This means that every \textit{rid} and \textit{step-order} uniquely identifies the description of the step, and subsequently the duration for such. Furthermore, as a step can involve multiple ingredients, the ingredients are not an identifying attribute of a step in a recipe.
    \item Every recipe step has at least one ingredient \\[0.5em]
    \textbf{Explanation:} The description states that a recipe step must include at least one ingredient. Subsequently, this becomes an intuitive requirement no inclusion of an ingredient does not change the current state of a recipe.
    \item \textit{amount} is determined by (\textit{rid, step-order, ingredient})\\[0.5em]
    \textbf{Explanation:} Different recipes may require the same ingredient, but different amounts. We identify the amount needed of an ingredient by the recipe, the step in which it is required, and the ingredient being used.
    \item \textit{ingredient} can have many \textit{UPC} \\[0.5em]
    \textbf{Explanation:} Many producers can produce the same ingredient. Therefore, ingredient can correspond to many different UPCs, that represent the offering of the ingredient by different producers.
    \item Every \textit{UPC} uniquely identifies a (\textit{ingredient, producer, quantity, price, unit-price})\\[0.5em]
    \textbf{Explanation:} A UPC uniquely identifies a product offering of an ingredient. As an ingredient can be produced by many producers, to differentiate the offering of such ingredients, we use the UPC to identify the producer, quantity, price, and unit-price of the ingredient.
    \item \textit{producer} can produce many \textit{UPC}\\[0.5em]
    \textbf{Explanation:} A producer can produce many different ingredients. Therefore, there can exist many different UPCs that correspond to different offerings of ingredients, with the same producer.
\end{itemize}

\subsubsection{}


\end{document}